% Part: first-order-logic
% Chapter: syntax-and-semantics
% Section: free-vars-sentences

\documentclass[../../../include/open-logic-section]{subfiles}

\begin{document}

\olfileid[hi]{fol}{syn}{fvs}

\olsection{मुक्त \printtoken{P}{variable} और \printtoken{P}{sentence} की धारणा}

\begin{defn}[!!a{variable} की मुक्त उपस्थितियाँ]
\ollabel{defn:free-occ}
किसी !!a{formula} में !!a{variable} की \emph{मुक्त} उपस्थितियाँ निम्न प्रकार
आगमनात्मक रूप से परिभाषित हैं:
\begin{enumerate}
\item \indcase*{!A}{$!A$ परमाण्विक है}{$\indfrm$ में !!{variable} की सभी
  उपस्थितियाँ मुक्त हैं।}

\tagitem{prvNot}{\indcase{!A}{\lnot !B}{$\indfrm$ में !!{variable} की मुक्त
  उपस्थितियाँ ठीक वही हैं जो $!B$ में हैं।}}{}

\item \indcase{!A}{(!B \ast !C)}{$\indfrm$ में !!{variable} की मुक्त
  उपस्थितियाँ वे हैं जो $!B$ में और वे जो~$!C$ में हैं।}

\tagitem{prvAll}{\indcase{!A}{\lforall[x][!B]}{$\indfrm$ में !!{variable}
  की मुक्त उपस्थितियाँ~$!B$ की सभी मुक्त उपस्थितियाँ हैं, सिवाय~$x$ की
  उपस्थितियों के।}}{}

\tagitem{prvEx}{\indcase{!A}{\lexists[x][!B]}{$\indfrm$ में !!{variable}
  की मुक्त उपस्थितियाँ~$!B$ की सभी मुक्त उपस्थितियाँ हैं, सिवाय~$x$ की
  उपस्थितियों के।}}{}
\end{enumerate}
\end{defn}

\begin{defn}[आबद्ध चर]
किसी सूत्र~$!A$ में !!a{variable} की कोई उपस्थिति \emph{आबद्ध} है यदि वह
मुक्त नहीं है।
\end{defn}

\begin{prob}
\olref[fol][syn][fvs]{defn:free-occ} की तर्ज पर चर की आबद्ध उपस्थितियों की
आगमनात्मक परिभाषा दीजिए।
\end{prob}

\begin{defn}[परास]
\iftag{prvAll}{यदि $\lforall[x][!B]$ किसी सूत्र~$!A$ में किसी उपसूत्र की
  उपस्थिति है, तो~$!B$ की संगत उपस्थिति को~$!A$ में~$\lforall[x]$ की संगत
  उपस्थिति का \emph{परास} कहते हैं। \iftag{prvEx}{$\lexists[x]$ के लिए भी
  इसी प्रकार।}{}}{यदि $\lexists[x][!B]$ किसी सूत्र~$!A$ में किसी उपसूत्र की
  उपस्थिति है, तो~$!B$ की संगत उपस्थिति को~$!A$ में~$\lexists[x]$ की संगत
  उपस्थिति का \emph{परास} कहते हैं।}

यदि $!B$ परिमाणक की किसी उपस्थिति
\iftag{prvAll}{$\lforall[x]$\iftag{prvEx}{ अथवा
    $\lexists[x]$}{}}{$\lexists[x]$} का सूत्र~$!A$ में परास है, तो~$x$ की~$!B$ में मुक्त
उपस्थितियाँ \iftag{prvAll}{$\lforall[x][!B]$\iftag{prvEx}{ और
$\lexists[x][!B]$}{}}{$\lexists[x][!B]$} में आबद्ध हैं। हम कहते हैं कि इन
उपस्थितियों को उक्त परिमाणक-उपस्थिति ने \emph{बाँधा} है।
\end{defn}

\begin{ex}
निम्न सूत्र पर विचार कीजिए:
\[
\lexists[\Obj v_0][\underbrace{\Atom{\Obj A^2_0}{\Obj v_0,\Obj v_1}}_{!B}] 
\]
$!B$, $\lexists[\Obj v_0]$ का परास दर्शाता है। परिमाणक $\Obj v_0$ की
उपस्थिति को $!B$ में बाँधता है, किंतु $\Obj v_1$ की उपस्थिति को नहीं बाँधता। अतः
इस स्थिति में $\Obj v_1$ मुक्त चर है।


अब हम देख सकते हैं कि किसी अधिक जटिल !!{formula}~$!A$ में यह कैसे काम करता
है:
\[
\lforall[\Obj v_0][\underbrace{(\Atom{\Obj A^1_0}{\Obj v_0} \lif
    \Atom{\Obj A^2_0}{\Obj v_0, \Obj v_1})}_{!B}] \lif \lexists[\Obj
  v_1][\underbrace{(\Atom{\Obj A^2_1}{\Obj v_0, \Obj v_1} \lor \lforall[\Obj v_0][\overbrace{\lnot \Atom{\Obj A^1_1}{\Obj v_0}}^{!D}])}_{!C}]
\]
$!B$ पहले $\lforall[\Obj v_0]$ का परास है, $!C$, $\lexists[\Obj v_1]$ का
परास है, और $!D$ दूसरे $\lforall[\Obj v_0]$ का परास है। पहला
$\lforall[\Obj v_0]$, $\Obj v_0$ की उपस्थितियों को~$!B$ में बाँधता है;
$\lexists[\Obj v_1]$, $\Obj v_1$ की उपस्थिति को $!C$ में बाँधता है; और
दूसरा $\lforall[\Obj v_0]$, $\Obj v_0$ की उपस्थिति को~$!D$ में बाँधता है।
$\Obj v_1$ की पहली उपस्थिति और $\Obj v_0$ की चौथी उपस्थिति~$!A$ में मुक्त
हैं। $\Obj v_0$ की अंतिम उपस्थिति $!D$ में मुक्त, किंतु $!C$ और~$!A$ में
आबद्ध है।
\end{ex}

\begin{defn}[वाक्य]
!!^a{formula}~$!A$ कोई \article{sentence} \emph{!!{sentence}} तभी और केवल
तभी है जब उसमें !!{variable}s की कोई मुक्त उपस्थिति न हो।
\end{defn}

% add examples!

\end{document}
